\documentclass[11pt]{article}
\usepackage[a4paper,margin=27mm]{geometry}
\usepackage{amsmath,amssymb,mathtools}
\usepackage{fontspec,unicode-math}
\setmainfont{Latin Modern Roman}
\setmathfont{Latin Modern Math}
\usepackage{xurl,hyperref}
\hypersetup{colorlinks=true,urlcolor=blue,linkcolor=blue}
\newcommand{\sref}[2]{\hyperlink{source:#1:#2}{[S#2]}}
\newcommand{\cataloguescope}[1]{\paragraph{Mathematical question.} #1}
\begin{document}
\section{Conjecture 3.19 — Rohrlich}\label{statement}

\subsection{Definitions and mathematical statement}

\paragraph{Shared notation.}$\mathbb N=\{1,2,\ldots\}$, and $\mathbb Z,\mathbb Q,\mathbb R,\mathbb C$ denote the integers, rationals, reals and complex numbers. Any other conventions are specified locally.


Let $G(z)=\Gamma(z)/\sqrt{2\pi}$ and let $H=\mathbb C^\times/\overline{\mathbb Q}^{\times}$, regarded as an abelian group. The functional equation of $\Gamma$ makes $g(x)=[G(x)]$ a well-defined map on $A=(\mathbb Q/\mathbb Z)\setminus\{0\}$. An odd distribution into an additive abelian group $B$ is a map $f:A\to B$ satisfying $f(-x)=-f(x)$ and $\sum_{j=0}^{N-1}f(x+j/N)=f(Nx)$ whenever $N\ge1$ and $Nx\ne0$ in $\mathbb Q/\mathbb Z$. Let $U$ be the free $\mathbb Z[1/2]$-module on symbols $[x]$ for $x\in A$, modulo exactly these oddness and distribution relations.
\paragraph{Statement recorded in the supplied source.}
Is the natural map $U\to H$, sending $[x]$ to $g(x)$, an isomorphism onto the subgroup generated by the $g(x)$ over $\mathbb Z[1/2]$? Equivalently, does every odd distribution with values in a group in which multiplication by $2$ is bijective factor uniquely through this gamma-value subgroup? \sref{statement}{2}
\subsection{Short English statement}
Are the gamma functional equations, reflection and multiplication relations the complete multiplicative relations among normalized gamma values at rational nonintegral points, after allowing division by two?
\subsection{Sources}
\begin{description}
\item[\hypertarget{source:statement:1}{S1}] ULAM.AI and the UnsolvedMath Contributors, UnsolvedMath: A Curated Collection of Open Mathematics Problems, catalogue entry. CC BY 4.0. \url{https://huggingface.co/datasets/ulamai/UnsolvedMath}.
\item[\hypertarget{source:statement:2}{S2}] Michel Waldschmidt, Open Diophantine Problems, arXiv:math/0312440v2 (24 January 2004), Conjecture 3.19 and its defining distribution equations. \url{https://arxiv.org/abs/math/0312440}.
\end{description}
\label{end:statement}
\end{document}
